%\generalExpl{When you do your literature study, you should have a nearly complete Chapters 1 and 2. You may also find it convenient to introduce the future work section into your report early – so that you can put things that you think about but decide not to do now into this section. Note that later you can move things between this future work section and what you have done as you may change your mind about what to do now versus what to put off to future work.}
%\generalExpl{What does a reader (another x student -- where x is your study line) need to know to understand your report? What have others already done? (This is the “related work”.) Explain what and how prior work/prior research will be applied on or used in the degree project/work (described in this thesis). Explain why and what is not used in the degree project and give valid reasons for rejecting the work/research.}

This chapter provides basic background information about xxx. Additionally, this chapter describes xxx. The chapter also describes related work xxxx.

\section{State Machine Replication}
A common way to provide fault tolerance and thereby improve availability for a service is to operate multiple servers that maintain an identical state. These replicas can serve as backups, shadowing a primary node, ready to take over during a failure, or be potentially globally distributed to handle requests from local clients. Many services, like databases, require a strict notion of consistency, leading to nodes having to coordinate their actions to prevent state divergence in both of these cases.

One approach is to decide on the execution of the same sequence of deterministic operations across all nodes. This is referred to as \gls{SMR} \cite{Schneider-ReplicatedStateMachines-1990}. The Raft protocol \cite{Ongaro-Raft-2014} popularized the concept of a replicated log as the central abstraction for developers of such replicated services. It is specifically designed around managing a linear, append-only log, which stores all \gls{SMR} operations for the application. Because this log is identical across all replicas, executing the operations in the order determined by it ensures that every node arrives at the same final state.

The underlying challenge addressed by Raft and similar protocols is uniform sequence consensus, which is detailed in Section \ref{sect:uniform_sequence_consensus}. Figure \ref{fig:state_machine_replication} illustrates how a replicated log serves as the backbone of a fault-tolerant system.

\begin{figure}[t]
    \centering
    \begin{tikzpicture}[
    node distance=1.0cm and 2.0cm,
    every node/.style={font=\small},
    component/.style={ rectangle, rounded corners=3pt, draw=black!70, thick, minimum width=4.8cm, minimum height=0.95cm, align=center },
    rsm/.style={component, fill=blue!25},
    log/.style={component, fill=orange!30},
    cons/.style={component, fill=green!25},
    client/.style={ rectangle, rounded corners=2pt, draw=black!80, thick, fill=white, minimum width=4cm, minimum height=0.9cm, font=\bfseries },
    serverFrame/.style={ rectangle, rounded corners=8pt, draw=black!80, very thick, inner sep=12pt },
    arrow/.style={ -{Latex[length=2.5mm]}, thick, draw=black!80 }
]
    % --- SERVER 1 ---
    \node[rsm] (rsm1) {Replicated State Machine};
    \node[log, below=of rsm1] (log1) {Replicated Log};
    \node[cons, below=of log1] (cons1) {Uniform Sequence Consensus};
    \node[serverFrame, fit=(rsm1) (cons1)] (frame1) {};
    \node[anchor=north, font=\bfseries, yshift=-5pt] at (frame1.south) {Server 1};
    
    % --- SERVER N ---
    \node[rsm, right=of rsm1, xshift=0.4cm] (rsmN) {Replicated State Machine};
    \node[log, below=of rsmN] (logN) {Replicated Log};
    \node[cons, below=of logN] (consN) {Uniform Sequence Consensus};
    \node[serverFrame, fit=(rsmN) (consN)] (frameN) {};
    \node[anchor=north, font=\bfseries, yshift=-5pt] at (frameN.south) {Server n};
    % --- CENTRAL ELLIPSIS --- 
    \node at ($(frame1.east)!0.5!(frameN.west)$) (dots) {\large\bfseries \dots};
    % --- CLIENTS ---
    \node[client, above=1.0cm of rsm1, xshift=-6pt, yshift=6pt] (cl1b) {};
    \node[client, above=1.0cm of rsm1, xshift=-3pt, yshift=3pt] (cl1c) {};
    \node[client, above=1.0cm of rsm1] (cl1) {Clients};
    \node[client, above=1.0cm of rsmN, xshift=-6pt, yshift=6pt] (clNb) {};
    \node[client, above=1.0cm of rsmN, xshift=-3pt, yshift=3pt] (clNc) {};
    \node[client, above=1.0cm of rsmN] (clN) {Clients};
    % --- ARROWS ---
    \foreach \i in {1, N} {
        % Vertical bidirectional flows
        \draw[-{Latex[length=2.5mm]}] ([xshift=-15pt]rsm\i.south) -- ([xshift=-15pt]log\i.north);
        \draw[-{Latex[length=2.5mm]}] ([xshift=15pt]log\i.north) -- ([xshift=15pt]rsm\i.south);
        \draw[-{Latex[length=2.5mm]}] ([xshift=-15pt]log\i.south) -- ([xshift=-15pt]cons\i.north);
        \draw[-{Latex[length=2.5mm]}] ([xshift=15pt]cons\i.north) -- ([xshift=15pt]log\i.south);
    }
    % Client Communication
    \draw[{Latex[length=2.5mm]}-{Latex[length=2.5mm]}] (cl1.south) -- (rsm1.north);
    \draw[{Latex[length=2.5mm]}-{Latex[length=2.5mm]}] (clN.south) -- (rsmN.north);
    % Inter-Server Consensus
    \draw[{Latex[length=2.5mm]}-{Latex[length=2.5mm]}] (cons1.east) -- (consN.west); 
\end{tikzpicture}
    \caption{Architecture of a \gls{SMR} service using a replicated log abstraction.}
    \label{fig:state_machine_replication}
\end{figure}


\subsection{Uniform Sequence Consensus}
\label{sect:uniform_sequence_consensus}
As described, the problem of agreeing on a value between coordinating processes lies at the core of fault-tolerant distributed systems. Any algorithm that relies on building a common state or synchronizing an action in an environment subject to partial failures must, in some form, solve the consensus problem.

We can formalize the abstraction layer of \textit{uniform consensus} according to Cachin et al. \cite{Cachin-ReliableDistributedProgramming-2011}. In this model, every process can propose a value to the consensus layer, which eventually outputs a final decided value. An algorithm for uniform consensus must satisfy the following properties:
\begin{enumerate}
  \item[] \textit{Validity:} Every decided value was previously proposed.
  \item[] \textit{Integrity:} No process decides more than once.
  \item[] \textit{Uniform agreement:} No two processes decide different values.
  \item[] \textit{Termination:} Every correct process eventually decides some value.
\end{enumerate}

Since uniform consensus is impossible to achieve in a purely \textit{asynchronous system} with no bounds on clock drift, transmission delay, and processing time \cite{Fischer-ConsensusAsynchronousSystem-1985}, we consider networks like the Internet as \textit{partially synchronous systems} throughout this thesis. This model assumes that while delays occur, there are sufficiently long stable periods for algorithms to make progress.

Standard uniform consensus only deals with a single decision, but a replicated log requires agreement on an ordered sequence. In an extension of this model, every participating process can propose values to the consensus layer, which eventually decides on a sequence of values. The relevant properties of this \textit{Uniform Sequence Consensus} have been defined by Haridi et al. \cite{Haridi-LeaderBasedSequencePaxos-2020}:
\begin{enumerate}
  \item[] \textit{Validity:} Every value in a sequence decided by a process was previously proposed.
  \item[] \textit{Integrity:} If a process decides a value sequence, it is a strict prefix of every future decided sequence by that server. 
  \item[] \textit{Uniform agreement:}  If two processes decide sequences of values, one sequence is a prefix of the other.
  \item[] \textit{Termination:} If a value is proposed infinitely often by a correct process, then every correct process eventually decides on a sequence containing this value.
\end{enumerate}

Ultimately, these properties provide the formal framework necessary to ensure that a distributed log remains a consistent single source of truth, allowing a collection of nodes to function as a unified, reliable service at the application level.

\subsection{Leader-Based Sequence Paxos}
Paxos denotes a family of algorithms designed to provide the guarantees of uniform (sequence) consensus. Its simplest form, often referred to as Basic Paxos or Single-Value Paxos, was originally introduced by Leslie Lamport \cite{Lamport-Paxos-2001} and solves uniform consensus within a partially synchronous system. It ensures that even under concurrent proposals from multiple nodes, the system converges on a single decision that satisfies the validity, integrity, and agreement properties defined in Section \ref{sect:uniform_sequence_consensus}. Understanding this underlying mechanism is essential before addressing the complexities of maintaining a continuous, ordered log.

Paxos decomposes the consensus process into three distinct phases. In the \textit{prepare phase}, a process acting as a \textit{proposer} attempts to claim the right to propose a value. It selects a unique, monotonically increasing proposal number $n$ and broadcasts a \textit{prepare} message containing $n$. Upon receipt, an \textit{acceptor} returns a \textit{promise} to the proposer, committing to reject any values from proposers that initiated a round with a lower proposal number than $n$. Crucially, if the acceptor has already accepted a value from a previous round, it includes that value and its associated proposal number in its response, forcing the proposer to respect any existing progress.

Once a proposer has secured promises from a majority (\textit{quorum}) of acceptors, it proceeds to the \textit{accept phase}, where it tries to make the acceptors lock in its proposal. It sends an \textit{accept} message to the promised acceptors, containing the proposal number $n$ and a value $v$ to be accepted. If the previous phase reported an existing value, the proposer must use the one with the highest proposal number; otherwise, it is free to use its own. An acceptor acknowledges this via an \textit{accepted} message, provided it has not already promised to a higher proposal number than $n$ previously.

As soon as a quorum of acceptors has accepted the same value, it is considered \textit{chosen}. At this stage, no other process can ever decide on a different value for this instance. Any participant that notices the value being chosen can notify the processes interested in the result (\textit{learners}). This is often done by the original proposer by broadcasting a \textit{decide} message to the learners. An example execution of a Paxos instance is displayed in Figure \ref{fig:paxos} with a total of three acceptors, who simultaneously act as proposers and learners.

\begin{figure}[t]
    \centering
    \begin{tikzpicture}[
    node distance=1.5cm and 2.5cm,
    every node/.style={font=\small},
    % Process styles
    process/.style={draw, diamond, inner sep=2pt, font=\bfseries\normalsize, minimum size=1.2cm, thick, fill=white},
    timeline/.style={thick, draw=black!90},
    % Arrow styles
    arrow/.style={-{Latex[length=2.5mm]}, thick},
    prepare/.style={arrow, draw=red!80!black},
    accept/.style={arrow, draw=green!60!black},
    learn/.style={arrow, draw=blue!70!black},
    % Helper styles
    phaseLine/.style={dashed, thick, draw=black!60},
    phaseLabel/.style={font=\bfseries\itshape},
    bottomLabel/.style={font=\small}
]

% --- PROCESS TIMELINES ---
\foreach \i in {1, 2, 3} {
    \node[process] (p\i) at (0.6, {-\i*1.8}) {P\i};
    \draw[timeline] (p\i.east) -- (12, {-\i*1.8});
}

% --- VERTICAL PHASE SEPARATORS ---
\draw[phaseLine] (4.8, -1.0) -- (4.8, -6.6);
\draw[phaseLine] (8.4, -1.0) -- (8.4, -6.6);

% --- PHASE 1: PREPARE ---
\node[phaseLabel, red!80!black] at (3.0, -0.75) {Prepare Phase};

% P1 self loop
\draw[prepare] (1.2, -1.8) .. controls (1.8, -1.1) and (2.4, -1.1) .. (3.0, -1.8);
% Messages to P2 and P3
\draw[prepare] (1.2, -1.8) -- (3.0, -3.6);
\draw[prepare] (1.2, -1.8) -- (3.0, -5.4);
% Responses back to P1
\draw[prepare] (3.0, -3.6) -- (4.8, -1.8);
\draw[prepare] (3.0, -5.4) -- (7.5, -1.8);

\node[bottomLabel] at (3.0, -6.2) {Prepare $|$ Promise};

% --- PHASE 2: ACCEPT ---
\node[phaseLabel, green!60!black] at (6.6, -0.75) {Accept Phase};

% P1 self loop
\draw[accept] (4.8, -1.8) .. controls (5.4, -1.1) and (6.0, -1.1) .. (6.6, -1.8);
% Messages to P2 and P3
\draw[accept] (4.8, -1.8) -- (6.6, -3.6);
\draw[accept] (4.8, -1.8) -- (6.6, -5.4);
% Responses back to P1
\draw[accept] (6.6, -3.6) -- (11.1, -1.8);
\draw[accept] (6.6, -5.4) -- (8.4, -1.8);

\node[bottomLabel] at (6.6, -6.2) {Accept $|$ Accepted};

% --- PHASE 3: LEARN ---
\node[phaseLabel, blue!70!black] at (10.2, -0.75) {Learn Phase};

% P1 self loop
\draw[learn] (8.4, -1.8) .. controls (9.0, -1.1) and (9.6, -1.1) .. (10.2, -1.8);
% Messages to P2 and P3
\draw[learn] (8.4, -1.8) -- (10.2, -3.6);
\draw[learn] (8.4, -1.8) -- (10.2, -5.4);

\node[bottomLabel] at (10.2, -6.2) {Decide};

\end{tikzpicture}
    \caption{Example execution of Paxos with three servers.}
    \label{fig:paxos}
\end{figure}

To satisfy the termination property of uniform consensus, a leader election mechanism can be used to elect a single distinguished proposer, who will eventually be the only process allowed to propose a value \cite{Lamport-Paxos-2001}. This principle leads directly to \textit{leader-based Sequence Paxos}, an extension of Paxos that handles an ordered sequence of values and provides the guarantees of uniform sequence consensus (see Section \ref{sect:uniform_sequence_consensus}).

The primary optimization of this approach is that if the leader remains stable, the prepare phase can be performed once for a whole series of log entries, rather than once per entry. This allows the system to effectively pipeline values through the accept phase using the elected leader as the distinguished proposer \cite{Lamport-Paxos-2001, Ng-OmniPaxos-2023, Haridi-LeaderBasedSequencePaxos-2020}. While specific implementations may augment the prepare phase with log-synchronization mechanisms depending on the leader election algorithm, this main idea stays the same across all implementations of uniform sequence consensus. Figure \ref{fig:leader_based_sequence_paxos} displays a common execution of a leader-based Sequence Paxos algorithm.


\begin{figure}[t]
    \centering
    \begin{tikzpicture}[
    node distance=1.5cm and 2.5cm,
    every node/.style={font=\small},
    % Process styles
    process/.style={draw, diamond, inner sep=2pt, font=\bfseries\normalsize, minimum size=1.0cm, thick, fill=white},
    timeline/.style={thick, draw=black!90},
    % Arrow styles
    arrow/.style={-{Latex[length=2.5mm]}, thick},
    prepare/.style={arrow, draw=red!80!black},
    accept1/.style={arrow, draw=green!60!black},
    accept2/.style={arrow, draw=green!100!black},
    accept3/.style={arrow, draw=green!40!black},
    % Helper styles
    phaseLine/.style={dashed, thick, draw=black!60},
    phaseLabel/.style={font=\bfseries\itshape},
    bottomLabel/.style={font=\small}
]

% --- PROCESS TIMELINES ---
\foreach \i in {1, 2, 3, 4, 5} {
    \node[process] (p\i) at (0.5, {-\i*1.2}) {P\i};
    \draw[timeline] (p\i.east) -- (12, {-\i*1.2});
}
\node[orange, yshift=0.2cm] at (p1.north) {\large\faCrown}; 

% --- VERTICAL PHASE SEPARATOR ---
\draw[phaseLine] (3.0, -0.0) -- (3.0, -6.8);

% --- PREPARE ---
\node[phaseLabel, red!80!black] at (2.0, -0.3) {Prepare};
% P1 self loop
\draw[prepare] (1.0, -1.2) .. controls (1.4, -0.6) and (1.6, -0.6) .. (2.0, -1.2);
% Messages to P2 and P3
\draw[prepare] (1.0, -1.2) -- (2.0, -2.4);
\draw[prepare] (1.0, -1.2) -- (2.0, -3.6);
% Responses back to P1
\draw[prepare] (2.0, -2.4) -- (3.0, -1.2);
\draw[prepare] (2.0, -3.6) -- (3.0, -1.2);

\node[phaseLabel, green!60!black] at (7.5, -0.3) {Accept (Pipelined)};
% --- ACCEPT 1 ---
% P1 self loop
\draw[accept1] (3.0, -1.2) .. controls (3.4, -0.6) and (3.6, -0.6) .. (4.0, -1.2);
% Messages to P2 and P3
\draw[accept1] (3.0, -1.2) -- (4.0, -2.4);
\draw[accept1] (3.0, -1.2) -- (4.0, -3.6);
% Responses back to P1
\draw[accept1] (4.0, -2.4) -- (8.5, -1.2);
\draw[accept1] (4.0, -3.6) -- (5.0, -1.2);

\draw[phaseLine] (8.5, -1.2) -- (8.5, -6.2);
\node[bottomLabel] at (8.5, -6.4) {Value 1};
\node[bottomLabel] at (8.5, -6.8) {Chosen};
% --- ACCEPT 2 ---
% P1 self loop
\draw[accept2] (4.5, -1.2) .. controls (4.9, -0.6) and (5.1, -0.6) .. (5.5, -1.2);
% Messages to P2 and P3
\draw[accept2] (4.5, -1.2) -- (6.5, -2.4);
\draw[accept2] (4.5, -1.2) -- (5.5, -3.6);
% Responses back to P1
\draw[accept2] (6.5, -2.4) -- (10.0, -1.2);
\draw[accept2] (5.5, -3.6) -- (6.5, -1.2);

\draw[phaseLine] (10.0, -1.2) -- (10.0, -6.2);
\node[bottomLabel] at (10.0, -6.4) {Value 2};
\node[bottomLabel] at (10.0, -6.8) {Chosen};
% --- ACCEPT 3 ---
% P1 self loop
\draw[accept3] (7.0, -1.2) .. controls (7.4, -0.6) and (7.6, -0.6) .. (8.0, -1.2);
% Messages to P2 and P3
\draw[accept3] (7.0, -1.2) -- (8.0, -2.4);
\draw[accept3] (7.0, -1.2) -- (8.0, -3.6);
% Responses back to P1
\draw[accept3] (8.0, -2.4) -- (10.5, -1.2);
\draw[accept3] (8.0, -3.6) -- (11.5, -1.2);

\draw[phaseLine] (11.5, -1.2) -- (11.5, -6.2);
\node[bottomLabel] at (11.5, -6.4) {Value 3};
\node[bottomLabel] at (11.5, -6.8) {Chosen};
\end{tikzpicture}
    \caption{Example execution of leader-based Sequence Paxos with five servers. Chosen values can be broadcast to learners similar to before.}
    \label{fig:leader_based_sequence_paxos}
\end{figure}

\subsection{Fast-Path Integration}
In the leader-based Sequence Paxos algorithm described above, only the leader can propose a value to be appended to the log. Thereby, there are a total of three message delays in a setting where every process might want to propose values \cite{Junqueira-ClassicVsFastPaxos-2007}:
\begin{enumerate}
    \item The original proposer forwards its value to the leader to start the append operation.
    \item The leader sends accept messages to all promised acceptors.
    \item The acceptors answer with corresponding accepted messages either to the learners directly or to the leader, who then informs the learners when a value is chosen. 
\end{enumerate}

The main idea behind the introduction of Fast Paxos is to save one of the communication steps by allowing any of the participating processes to directly propose values to the acceptors instead of first forwarding them to the leader. This was also originally devised by Leslie Lamport \cite{Lamport-FastPaxos-2006}. In such a fast round, the acceptors now simply accept the first value they receive from any proposer and notify the leader (and possibly the learners) of their decision using an accepted message. A typical flow of messages with successful fast paths is depicted in Figure \ref{fig:leader_based_sequence_paxos_fast_path}.

\begin{figure}[t]
    \centering
    \begin{tikzpicture}[
    node distance=1.5cm and 2.5cm,
    every node/.style={font=\small},
    % Process styles
    process/.style={draw, diamond, inner sep=2pt, font=\bfseries\normalsize, minimum size=1.0cm, thick, fill=white},
    timeline/.style={thick, draw=black!90},
    % Arrow styles
    arrow/.style={-{Latex[length=2.5mm]}, thick},
    prepare/.style={arrow, draw=red!80!black},
    accept1/.style={arrow, draw=green!60!black},
    accept2/.style={arrow, draw=green!100!black},
    accept3/.style={arrow, draw=green!40!black},
    % Helper styles
    phaseLine/.style={dashed, thick, draw=black!60},
    phaseLabel/.style={font=\bfseries\itshape},
    bottomLabel/.style={font=\small}
]

% --- PROCESS TIMELINES ---
\foreach \i in {1, 2, 3, 4, 5} {
    \node[process] (p\i) at (0.5, {-\i*1.2}) {P\i};
    \draw[timeline] (p\i.east) -- (12, {-\i*1.2});
}
\node[orange, yshift=0.2cm] at (p1.north) {\large\faCrown}; 

% --- VERTICAL PHASE SEPARATOR ---
\draw[phaseLine] (3.0, -0.0) -- (3.0, -6.8);

% --- PREPARE ---
\node[phaseLabel, red!80!black] at (2.0, -0.3) {Prepare};
% P1 self loop
\draw[prepare] (1.0, -1.2) .. controls (1.4, -0.6) and (1.6, -0.6) .. (2.0, -1.2);
% Messages to P2, P3, P4
\draw[prepare] (1.0, -1.2) -- (2.0, -2.4);
\draw[prepare] (1.0, -1.2) -- (2.0, -3.6);
\draw[prepare] (1.0, -1.2) -- (2.0, -4.8);
% Responses back to P1
\draw[prepare] (2.0, -2.4) -- (3.0, -1.2);
\draw[prepare] (2.0, -3.6) -- (3.0, -1.2);
\draw[prepare] (2.0, -4.8) -- (3.0, -1.2);

\node[phaseLabel, green!60!black] at (7.5, -0.3) {Accept (Pipelined)};

% --- ACCEPT 1 ---
% P1 self loop
\draw[accept1] (4.0, -1.2) .. controls (4.4, -0.6) and (4.6, -0.6) .. (5.0, -1.2);
% Messages to P1, P2 and P4
\draw[accept1] (3.0, -3.6) -- (4.0, -1.2);
\draw[accept1] (3.0, -3.6) -- (4.5, -2.4);
\draw[accept1] (3.0, -3.6) -- (4.0, -4.8);
% Responses to leader P1
\draw[accept1] (4.5, -2.4) -- (5.5, -1.2);
\draw[accept1] (4.0, -4.8) -- (6.0, -1.2);

\draw[phaseLine] (6.0, -1.2) -- (6.0, -6.2);
\node[bottomLabel] at (6.0, -6.4) {Value 1};
\node[bottomLabel] at (6.0, -6.8) {Chosen};

% --- ACCEPT 2 ---
% P1 self loop
\draw[accept2] (5.5, -1.2) .. controls (5.9, -0.6) and (6.1, -0.6) .. (6.5, -1.2);
% Messages to P2, P3 and P4
\draw[accept2] (5.5, -1.2) -- (8.0, -2.4);
\draw[accept2] (5.5, -1.2) -- (6.5, -3.6);
\draw[accept2] (5.5, -1.2) -- (6.5, -4.8);
% Responses to leader P1
\draw[accept2] (8.0, -2.4) -- (9.5, -1.2);
\draw[accept2] (6.5, -3.6) -- (7.5, -1.2);
\draw[accept2] (6.5, -4.8) -- (7.5, -1.2);

\draw[phaseLine] (9.5, -1.2) -- (9.5, -6.2);
\node[bottomLabel] at (9.5, -6.4) {Value 2};
\node[bottomLabel] at (9.5, -6.8) {Chosen};

% --- ACCEPT 3 ---
% P1 self loop
\draw[accept3] (9.0, -1.2) .. controls (9.4, -0.6) and (9.6, -0.6) .. (10.0, -1.2);
% Messages to P1, P2 and P3
\draw[accept3] (8.0, -4.8) -- (9.0, -1.2);
\draw[accept3] (8.0, -4.8) -- (9.0, -2.4);
\draw[accept3] (8.0, -4.8) -- (9.0, -3.6);
% Responses to leader P1
\draw[accept3] (9.0, -2.4) -- (11.0, -1.2);
\draw[accept3] (9.0, -3.6) -- (11.5, -1.2);

\draw[phaseLine] (11.5, -1.2) -- (11.5, -6.2);
\node[bottomLabel] at (11.5, -6.4) {Value 3};
\node[bottomLabel] at (11.5, -6.8) {Chosen};
\end{tikzpicture}
    \caption{Example execution of leader-based Sequence Paxos with five servers using the fast path.}
    \label{fig:leader_based_sequence_paxos_fast_path}
\end{figure}

Compared to the original algorithm discussed earlier, a simple majority quorum is not sufficient anymore to consider a value chosen when using the fast path. To guarantee safety, a quorum for a fast round (\textit{fast quorum}) must intersect with any majority of nodes in at least half of the nodes included in the majority \cite{Zhao-FastPaxos-2015}. In practice, the size of a fast quorum must be at least $\lceil \frac{3n}{4} \rceil$ for a system of $n$ acceptors. This also means that the fast-path algorithm can tolerate fewer failures in a round compared to the earlier design. 

Additionally, it is obvious that the fast path is optimistic. A collision can occur when multiple proposers send different values simultaneously, such that no single value achieves the fast quorum threshold. The leader can detect such a collision and initiate a regular Paxos round using the most accepted value, as seen in Figure \ref{fig:leader_based_sequence_paxos_fast_path_collision}. This \textit{slow path} requires two additional communication steps, which means that a fast round with a collision is even slower than a round following the original algorithm.

\begin{figure}[t]
    \centering
    \begin{tikzpicture}[
    node distance=1.5cm and 2.5cm,
    every node/.style={font=\small},
    % Process styles
    process/.style={draw, diamond, inner sep=2pt, font=\bfseries\normalsize, minimum size=1.0cm, thick, fill=white},
    timeline/.style={thick, draw=black!90},
    % Arrow styles
    arrow/.style={-{Latex[length=2.5mm]}, thick},
    accept1/.style={arrow, draw=green!60!black},
    accept2/.style={arrow, draw=green!100!black},
    accept3/.style={arrow, draw=green!40!black},
    % Helper styles
    phaseLine/.style={dashed, thick, draw=black!60},
    phaseLabel/.style={font=\bfseries\itshape},
    bottomLabel/.style={font=\small}
]

% --- PROCESS TIMELINES ---
\foreach \i in {1, 2, 3, 4, 5} {
    \node[process] (p\i) at (0.5, {-\i*1.2}) {P\i};
    \draw[timeline] (p\i.east) -- (12, {-\i*1.2});
}
\node[orange, yshift=0.2cm] at (p1.north) {\large\faCrown}; 

% --- ACCEPT 1 ---
% P1 self loop
\draw[accept1] (2.0, -1.2) .. controls (2.4, -0.6) and (2.6, -0.6) .. (3.0, -1.2);
% Messages to P1, P2, P4 and P5
\draw[accept1] (1.0, -3.6) -- (2.0, -1.2);
\draw[accept1] (1.0, -3.6) -- (2.0, -2.4);
\draw[accept1] (1.0, -3.6) -- (2.0, -4.8);
\draw[accept1] (1.0, -3.6) -- (2.0, -6.0);
% Responses to leader P1
\draw[accept1] (2.0, -2.4) -- (3.0, -1.2);
\draw[accept1] (2.0, -4.8) -- (3.0, -1.2);
\draw[accept1] (2.0, -6.0) -- (4.0, -1.2);

\draw[phaseLine] (3.0, -1.2) -- (3.0, -6.2);
\node[bottomLabel] at (3.0, -6.4) {Value 1};
\node[bottomLabel] at (3.0, -6.8) {Chosen};

% --- COLLISION ---
% P1 self loop
\draw[accept2] (5.5, -1.2) .. controls (5.9, -0.6) and (6.1, -0.6) .. (6.5, -1.2);
% Messages to P1, P3, P4
\draw[accept2] (4.5, -2.4) -- (5.5, -1.2);
\draw[accept2] (4.5, -2.4) -- (5.5, -3.6);
\draw[accept2] (4.5, -2.4) -- (6.0, -4.8);
% Responses to leader P1
\draw[accept2] (5.5, -3.6) -- (6.5, -1.2);

% Messages to P1, P2, P3, P5
\draw[accept3] (5.5, -4.8) -- (7.0, -1.2);
\draw[accept3] (5.5, -4.8) -- (7.0, -2.4);
\draw[accept3] (5.5, -4.8) -- (6.5, -3.6);
\draw[accept3] (5.5, -4.8) -- (6.5, -6.0);
% Responses to leader P1
\draw[accept3] (6.5, -6.0) -- (8.0, -1.2);

\draw[phaseLine] (8.0, -0.0) -- (8.0, -6.2);
\node[bottomLabel] at (8.0, -6.4) {Collision};
\node[bottomLabel] at (8.0, -6.8) {Detected};

\node[phaseLabel, green!60!black] at (5.5, -0.3) {Fast Path};

% --- RECOVERY ---
% P1 self loop
\draw[accept2] (8.0, -1.2) .. controls (8.4, -0.6) and (8.6, -0.6) .. (9.0, -1.2);
% Messages to P2, P3, P4
\draw[accept2] (8.0, -1.2) -- (9.0, -2.4);
\draw[accept2] (8.0, -1.2) -- (9.0, -3.6);
\draw[accept2] (8.0, -1.2) -- (9.0, -4.8);
\draw[accept2] (8.0, -1.2) -- (9.0, -6.0);
% Responses back to P1
\draw[accept2] (9.0, -2.4) -- (10.0, -1.2);
\draw[accept2] (9.0, -3.6) -- (10.0, -1.2);
\draw[accept2] (9.0, -4.8) -- (11.0, -1.2); 
\draw[accept2] (9.0, -6.0) -- (11.5, -1.2); 

\draw[phaseLine] (10.0, -1.2) -- (10.0, -6.2);
\node[bottomLabel] at (10.0, -6.4) {Value 2};
\node[bottomLabel] at (10.0, -6.8) {Chosen};

\node[phaseLabel, green!60!black] at (10.0, -0.3) {Slow Path};

\end{tikzpicture}
    \caption{Example execution of leader-based Sequence Paxos with five servers and a collision of proposals.}
    \label{fig:leader_based_sequence_paxos_fast_path_collision}
\end{figure}

\todo{polish}

\section{Conformal Prediction}
When deploying predictive models in high-stakes applications, understanding the uncertainty of its output is just as critical as assessing its actual accuracy on prediction tasks. This makes it \gls{UQ} a critical research field in statistics and machine learning. \gls{CP} describes a statistical framework for \gls{UQ} with the primary objective of transforming a heuristic notion of uncertainty into statistically rigorous prediction sets/intervals that are guaranteed to contain the ground truth with a user-defined probability. \gls{CP} has several desirable properties that make it suited for model deployment \cite{Shafer-ConformalPrediction-2008, Angelopoulus-ConformalPrediction-2023, Zhou-ConformalPredictionDataPerspective-2025}:
\begin{itemize}
    \item \textit{Model-agnostic:} No knowledge of the underlying model architecture is required when applying this method. 
    \item \textit{Distribution-free:} Valid coverage is provided under the assumption of data exchangeability. A sequence of random variables is considered exchangeable if their joint probability distribution is invariant under any permutation of this sequence. This assumption implies that the order in which data points are observed does not carry information about their underlying distribution.
    \item \textit{Efficient:} Because it functions as a lightweight wrapper, it can be applied to any existing model without the need for retraining.
    \item \textit{Informative:} It produces prediction sets/intervals that are narrow enough to provide actionable insights about the model's output.
    \item \textit{Finite-sample valid:} Coverage is guaranteed even if prediction is constructed based of a finite number of data points. 
\end{itemize}

The following sections introduce specific conformal methods and describe how they can be used for quantifying uncertainty in classification tasks.
\todo{Polish}

\subsection{Full Conformal Prediction}
\label{sect:full_conformal_prediction}
Given a training set $\mathcal{D}_n \in (\mathcal{X} \times \mathcal{Y})^n$ containing features and labels $(X_1, Y_1), (X_2, Y_2), ..., (X_n, Y_n)$ and a test point $(X_{test}, Y_{test}) \in \mathcal{X} \times \mathcal{Y}$ with the test label being unobserved, the goal of conformal prediction for classification is to construct a calibration set for the test feature $\mathcal{C}(X_{test}) \subseteq \mathcal{Y}$ such that it satisfies the \textit{marginal predictive coverage} for a user-defined error rate $\alpha$ \cite{Angelopoulus-ConformalPrediction-2023, Zhou-ConformalPredictionDataPerspective-2025, Angelopoulos-ConformalPredictionTheory-2025}:
$$\mathbb{P}(Y_{test} \in \mathcal{C}(X_{test})) \ge 1 - \alpha$$

This is done based on a \textit{non-conformity score} function $s((x, y); \mathcal{D}) \in \mathbb{R}$ that measures the error of a model trained with $\mathcal{D} \in (\mathcal{X} \times \mathcal{Y})^k$ on the feature-label pair $(x,y) \in \mathcal{X} \times \mathcal{Y}$. A higher score indicates a worse agreement between $x$ and $y$ based on the training data \cite{Angelopoulus-ConformalPrediction-2023}. The function itself is user-defined and primarily determines the usefulness of the resulting prediction sets, meaning it should be based off a heuristic notion of uncertainty.

The core idea of \textit{full \gls{CP}} is to compare the non-conformity scores of all hypothesized labels $y \in \mathcal{Y}$ for the given test feature $X_{test}$ with the scores of the training data. A comparably large score would hereby indicate that a particular label is inconsistent with the data, and the final prediction set $\mathcal{C}(X_{test})$ then simply contains all consistent labels. The consistency threshold is the $(1 - \alpha)$ quantile with some finite-sample correction, which satisfies the marginal predictive coverage for the error rate $\alpha$ \cite{Angelopoulos-ConformalPredictionTheory-2025}. Listing \ref{list:full_conformal_prediction} shows the detailed algorithm of full conformal prediction as adapted from Angelopoulos et al. \cite{Angelopoulos-ConformalPredictionTheory-2025}.

\begin{listing}[t]
    \begin{algorithmic}[1]
        \renewcommand{\algorithmicrequire}{}
        \REQUIRE
        \noindent\begin{tabular}{@{}ll}
            Training data: & $\mathcal{D} = \{(X_1, Y_1), \dots, (X_n, Y_n)\}\in (\mathcal{X} \times \mathcal{Y})^n$\\
            Test feature: & $X_{test} \in \mathcal{X}$\\
            Target error rate: & $\alpha \in [0\dots1]$\\
            Non-conformity score: & $s: (\mathcal{X} \times \mathcal{Y}) \times (\mathcal{X} \times \mathcal{Y})^k \rightarrow \mathbb{R}$
        \end{tabular}
        \vspace{2mm}
        \hrule
        \vspace{2mm}
        \FORALL{$y \in \mathcal{Y}$}
            \STATE Full training set: $\mathcal{D}^y \leftarrow \mathcal{D} \cup (X_{test}, y)$
            \STATE Conformal scores: $S_i^y \leftarrow s((X_i, Y_i); \mathcal{D}^y)$ for all $i \in \{1, \dots, n\}$
            \STATE Conformal scores of the hypothesized test point: $S_{test}^y \leftarrow s((X_{test}, y); \mathcal{D}^y)$
            \STATE Conformal quantile: $\hat{q}^y \leftarrow Quantile(S_1^y, \dots, S_n^y; (1 - \alpha)(1 + 1/n))$
        \ENDFOR
        \RETURN Prediction set: $\mathcal{C}(X_{test}) \leftarrow \{y \in \mathcal{Y} : S_{n+1}^y \le \hat{q}^y\}$
    \end{algorithmic}
    \caption{Full conformal prediction algorithm}
    \label{list:full_conformal_prediction}
\end{listing}

\subsection{Split Conformal Prediction}
\label{sect:split_conformal_prediction}

Full \gls{CP} as described is hugely inefficient. It requires full retraining for each possible label and a repetition of this process for every new observation. A computationally efficient variant called \textit{split \gls{CP}} instead divides the training set into data used for model fitting and data used as a calibration set for computing the non-conformity scores as a comparison \cite{Vovk-AlgorithmicLearning-2022}.

\begin{figure}[t]
    \centering
    \begin{tikzpicture}[
    every node/.style={font=\small},
    dataPoint/.style={circle, draw=black, thick, minimum size=8pt, inner sep=0pt},
    consistentData/.style={dataPoint, fill=green!40},
    inconsistentData/.style={dataPoint, fill=red!40},
    testData/.style={dataPoint, fill=blue!40},
    arrow/.style={-{Latex[length=2.5mm]}, thin, draw=gray},
    brace/.style={decorate, decoration={brace, amplitude=6pt}, thick, draw=black!70},
]
    
    % --- Main Horizontal Line ---
    \draw[black!70, thick] (0,0) -- (12,0);
    
    % --- Non-conformity Scores (Green Dots) ---
    \node[consistentData] (cp1) at (1.5,0) {};
    \node[consistentData] (cp2) at (2,0) {};
    \node[consistentData] (cp3) at (3.5,0) {};
    \node[consistentData] (cp4) at (4.2,0) {};
    \node[consistentData] (cp5) at (5,0) {};
    \node[consistentData] (cp6) at (6,0) {};
    \node[consistentData] (cp7) at (6.5,0) {};
    \node[consistentData] (cp8) at (8,0) {};
    
    % --- "Consistent" Label ---
    \node at (4,-1) {\textbf{consistent}};
    
    % --- Threshold Line ---
    \draw[black!70, thick, dashed] (8,2) -- (8, -0.5);
    \node[below=0.3 of cp8, anchor=north, font=\normalsize] (qHatLabel) {$\hat{q}$};
    
    % --- Inconsistent Region and Data Points ---
    \node[inconsistentData] (icp1) at (8.5,0) {};
    \node[inconsistentData] (icp2) at (10,0) {};
    \node at (10,-1) {\textbf{inconsistent}};
    
    % --- Braces ---
    \node (p1Anchor) at (0,0) {};
    \node[above=12pt of $(p1Anchor)!0.5!(cp1)$] {$1/(n+1)$}; 
    \draw[brace] (p1Anchor.north) -- (cp1.north);
    \node[above=12pt of $(cp3)!0.5!(cp4)$] {$1/(n+1)$}; 
    \draw[brace] (cp3.north) -- (cp4.north);
    \node[above=12pt of $(icp1)!0.5!(icp2)$] {$1/(n+1)$}; 
    \draw[brace] (icp1.north) -- (icp2.north);
    
    % --- Test Score ---
    \node[testData] (testScore) at (6, 2.0) {};
    \node[above=3pt of testScore, font=\normalsize] {$s(X_{test}, Y_{test})$};
    
    % --- Midpoints ---
    \path (p1Anchor.center) -- (cp1.center) coordinate[pos=0.5] (mid1);
    \path (cp3.center) -- (cp4.center) coordinate[pos=0.5] (mid2);
    \path (icp1.center) -- (icp2.center) coordinate[pos=0.5] (mid3);
    
    % --- Arrows ---
    \draw[arrow] (testScore.west) -- ($(mid1)+(0,0.30)$);
    \draw[arrow] (testScore.south) -- ($(mid2)+(0,0.30)$);
    \draw[arrow] (testScore.east) -- ($(mid3)+(0,0.30)$);
\end{tikzpicture}
    \caption{Example execution of leader-based Sequence Paxos with five servers and a collision of proposals.}
    \label{fig:split_conformal_prediction}
\end{figure}

\begin{listing}[t]
    \begin{algorithmic}[1]
        \renewcommand{\algorithmicrequire}{}
        \REQUIRE
        \noindent\begin{tabular}{@{}ll}
            Model fitting data: & $\mathcal{D_{mod}} \in (\mathcal{X} \times \mathcal{Y})^m$\\
            Calibration data: & $\mathcal{D_{cal}} = \{(X_1, Y_1), \dots, (X_n, Y_n)\}\in (\mathcal{X} \times \mathcal{Y})^n$\\
            Test feature: & $X_{test} \in \mathcal{X}$\\
            Target error rate: & $\alpha \in [0\dots1]$\\
        \end{tabular}
        \vspace{2mm}
        \hrule
        \vspace{2mm}
        \STATE Construct a score function $s: \mathcal{X} \times \mathcal{Y} \rightarrow \mathbb{R}$ using $\mathcal{D}_{mod}$
        \STATE Conformal scores of calibration set: $S_i \leftarrow s(X_i, Y_i)$, for all $i \in \{1, \dots, n\}$
        \STATE Conformal quantile: $\hat{q} \leftarrow Quantile(S_1, \dots, S_n; (1 - \alpha)(1 + 1/n))$
        \RETURN Prediction set: $\mathcal{C}(X_{test}) \leftarrow \{y \in \mathcal{Y}: s(X_{test}, y) \le \hat{q}\}$
    \end{algorithmic}
    \caption{Split conformal prediction algorithm}
    \label{list:split_conformal_prediction}
\end{listing}

exchangeability ensures that the "non-conformity score" of a previously unseen test observation is equally likely to occupy any rank when compared against a set of calibration scores.

\todo{Write out, make figure for split conformal showing the equal probability that a test point lands in between two calibration points}

\subsection{Conformal Risk Control}
\label{sect:conformal_risk_control}

\subsection{Adaptive Conformal Inference?}

\section{Related work area}

\subsection{Major related work 1}

\subsection{Major related work n}

\subsection{Minor related work 1}

\subsection{Minor related work n}


\section{Summary}
\sweExpl{Det är trevligt om detta kapitel avslutas med en sammanfattning. Till exempel kan du inkludera en tabell som sammanfattar andras idéer och fördelar och nackdelar med varje - så som senare kan du jämföra din lösning till var och en av dessa. Detta kommer också att hjälpa dig att definiera de variabler som du kommer att använda för din utvärdering.}

\engExpl{It is nice to have this chapter conclude with a summary. For example, you can include a table that summarizes other people's ideas and benefits and drawbacks with each - so as later you can compare your solution to each of them. This will also help you define the variables that you will use for your evaluation.}